\documentclass[11pt,a4paper]{article}
\usepackage[T1]{fontenc}
\usepackage[utf8]{inputenc}
\usepackage{amsmath,amssymb,amsthm}
\usepackage[margin=1in]{geometry}
\usepackage[colorlinks=true,linkcolor=blue,urlcolor=blue]{hyperref}
\theoremstyle{plain}
\newtheorem{theorem}{Theorem}
\newtheorem{lemma}[theorem]{Lemma}
\newtheorem{proposition}[theorem]{Proposition}
\newtheorem{corollary}[theorem]{Corollary}
\theoremstyle{definition}
\newtheorem{remark}[theorem]{Remark}

\title{The empirical recurrence for OEIS A210057, proved by counting patterns}
\author{Adrian Perez Fontelles\\ \small Independent researcher}
\date{31 August 2026}
\begin{document}
\maketitle

\begin{abstract}
OEIS A210057 counts arrays subject to a local condition \emph{and} to the clause ``new values
$0..3$ introduced in row major order'', which says the array is written in canonical
form. That clause is not local, so a transfer matrix over the alphabet does not see it. It
does not have to: what is being counted is not arrays but equality patterns, and the number
of patterns is recovered from the numbers $L_i$ of arrays over an $i$-letter alphabet by
inverting a falling factorial. The result is the clean identity
$4!\,a(n)=\sum_i\binom{4}{i}D_{4-i}L_i(n)$ with $D_m$ the derangement numbers,
and each $L_i$ is an ordinary walk count. A second reduction --- the chain is strongly
lumpable over the patterns, because the condition is relabelling-invariant --- brings the
state space down from $1300$ to $S=109$. The entry's empirical recurrence of order
$28$, contributed by R.~H.~Hardin, then follows from a finite exact computation.
\end{abstract}

\noindent\small 2020 Mathematics Subject Classification. 05A15, 05A18, 15B36.\normalsize

\section{The sequence and the conjecture}

OEIS A210057 is ``Number of (n+1)X5 0..3 arrays with every 2X2 subblock having one, two or four distinct values, and new values 0..3 introduced in row major order.''. It has offset $1$ and begins
\[
2187, 124097, 7293673, 438056462, 26586250017, 1622061565898, 99225470062505, 6077972350575098,\ \dots
\]
The entry carries the following comment:
\begin{quote}\small
Empirical: a(n) = 176*a(n-1) -12383*a(n-2) +457521*a(n-3) -9583668*a(n-4) +109074306*a(n-5) -432638439*a(n-6) -4419650349*a(n-7) +64230480793*a(n-8) -252540170918*a(n-9) -595255299234*a(n-10) +8433913817808*a(n-11) -21427208559976*a(n-12) -38069222537892*a(n-13) +289460472761408*a(n-14) -336643109299632*a(n-15) -883702002797616*a(n-16) +2460214504431920*a(n-17) -361513807179872*a(n-18) -4377499678974656*a(n-19) +3253938371960768*a(n-20) +2790969020407040*a(n-21) -3141539845016064*a(n-22) -669537234834176*a(n-23) +1122282161785600*a(n-24) +46793169024000*a(n-25) -143518636118016*a(n-26) -1001036759040*a(n-27) +3148485230592*a(n-28)
\end{quote}
As of the ``Last modified'' line on the live entry (Oct 03 2025 15:40:05, revision 10) the statement is
still recorded as empirical, and nothing on the entry records it as settled either way.

Written out, the claim is
\begin{equation}\label{eq:conj}
a(n)\;=\;176\,a(n-1) - 12383\,a(n-2) + 457521\,a(n-3) - 9583668\,a(n-4) + 109074306\,a(n-5) - 432638439\,a(n-6) - 4419650349\,a(n-7) + 64230480793\,a(n-8) - 252540170918\,a(n-9) - 595255299234\,a(n-10) + 8433913817808\,a(n-11) - 21427208559976\,a(n-12) - 38069222537892\,a(n-13) + 289460472761408\,a(n-14) - 336643109299632\,a(n-15) - 883702002797616\,a(n-16) + 2460214504431920\,a(n-17) - 361513807179872\,a(n-18) - 4377499678974656\,a(n-19) + 3253938371960768\,a(n-20) + 2790969020407040\,a(n-21) - 3141539845016064\,a(n-22) - 669537234834176\,a(n-23) + 1122282161785600\,a(n-24) + 46793169024000\,a(n-25) - 143518636118016\,a(n-26) - 1001036759040\,a(n-27) + 3148485230592\,a(n-28)
\end{equation}
for all $n>27$.

\section{What the canonical-form clause counts}

Fix the shape: $5$ columns and $L=n+1$ rows, so $L\cdot5$ cells. Call a
partition of the cells into nonempty classes a \emph{pattern}; an array over an alphabet
determines a pattern, two cells lying in the same class exactly when they carry the same
value.

Every $2\times2$ subblock of the array occupies two consecutive
rows and two consecutive columns; writing those two rows as $r$ and $s$ the
subblock is
\[
\begin{pmatrix}\alpha&\beta\\\gamma&\delta\end{pmatrix}=\begin{pmatrix}r_j&r_{j+1}\\s_j&s_{j+1}\end{pmatrix},
\]
and the entry's requirement on it is
\begin{equation}\label{eq:loc}
\#\{\alpha,\beta,\gamma,\delta\}\in\{1,2,4\}.
\end{equation}

The condition is stated purely in terms of equalities between entries, so it is invariant
under any relabelling of the values and is therefore a property of the pattern alone. Call a
pattern \emph{admissible} when it has the property.

The clause ``new values $0..3$ introduced in row major order'' selects, from each class
of arrays that differ only by a relabelling, exactly one representative: reading the cells in
row major order, the first occurrence of the value $v$ must precede the first occurrence of
$v+1$. Every pattern with at most $K=4$ classes has exactly one such representative, and
every array in canonical form is the representative of its own pattern. Hence, writing
$N_j(n)$ for the number of admissible patterns with exactly $j$ classes,
\begin{equation}\label{eq:apat}
a(n)\;=\;\sum_{j=0}^{4}N_j(n).
\end{equation}

\begin{lemma}\label{lem:inv}
Let $L_i(n)$ be the number of arrays over an alphabet of $i$ letters satisfying the local
condition, with no canonical-form clause. Then
\[
L_i(n)=\sum_j N_j(n)\,i(i-1)\cdots(i-j+1),
\qquad
4!\,\sum_{j=0}^{4}N_j(n)=\sum_{i=0}^{4}\binom{4}{i}D_{4-i}\,L_i(n),
\]
where $D_m=m!\sum_{t=0}^{m}(-1)^t/t!$ is the number of derangements of $m$ letters.
\end{lemma}

\begin{proof}
An array over $i$ letters satisfying the condition determines an admissible pattern together
with an injection of its classes into the alphabet, and conversely; a pattern with $j$
classes admits $i(i-1)\cdots(i-j+1)$ injections. That is the first identity. Writing
$i(i-1)\cdots(i-j+1)=j!\binom{i}{j}$ it reads $L_i=\sum_j (j!N_j)\binom{i}{j}$, so by
binomial inversion $j!N_j=\sum_{i}(-1)^{j-i}\binom{j}{i}L_i$. Summing over
$j\le 4$ and exchanging the order of summation,
\[
\sum_{j=0}^{4}N_j=\sum_{i=0}^{4}L_i\sum_{j=i}^{4}\frac{(-1)^{j-i}}{i!\,(j-i)!}
=\sum_{i=0}^{4}\frac{L_i}{i!}\sum_{t=0}^{4-i}\frac{(-1)^t}{t!} ,
\]
and multiplying by $4!$ turns the coefficient of $L_i$ into
$\frac{4!}{i!}\sum_{t\le 4-i}(-1)^t/t!=\binom{4}{i}(4-i)!\sum_{t\le 4-i}(-1)^t/t!
=\binom{4}{i}D_{4-i}$, an integer.
\end{proof}

\section{Each $L_i$ is a walk count, and the chain lumps}

Fix $i$ and let $V_i=\{0,\dots,i-1\}^{5}$ be the possible rows. For $r,s\in V_i$
declare $r\to s$ an edge when placing $s$ directly after $r$ satisfies \eqref{eq:loc} at
every column, and let $M_i$ be the adjacency matrix. An array is exactly a sequence of
rows, and its condition is exactly the conjunction of the edge conditions along that
sequence, so
\[
L_i(n)\;=\;\mathbf{s}_i^{\!\top}M_i^{\,n}\mathbf{1} .
\]
Because the condition involves no row before the first, every
row may start a walk, and $\mathbf{s}_i=\mathbf{1}$.

\begin{lemma}\label{lem:lump}
Let $\pi(r)$ be the equality pattern of the row $r$. For all $r,r'$ with
$\pi(r)=\pi(r')$ and every pattern $Q$,
\[
\#\{s:\ r\to s,\ \pi(s)=Q\}\;=\;\#\{s:\ r'\to s,\ \pi(s)=Q\}\, .
\]
Consequently the numbers $\widehat M_i[P][Q]$ obtained by taking any representative of $P$
are well defined, and with $e_P=i(i-1)\cdots(i-|P|+1)$ the number of rows of pattern $P$,
\[
\mathbf{s}_i^{\!\top}M_i^{\,m}\mathbf{1}\;=\;\sum_P \varepsilon_P\,f_m(P),
\qquad f_0\equiv1,\quad f_{m+1}(P)=\sum_Q\widehat M_i[P][Q]f_m(Q),
\]
where $\varepsilon_P=e_P$ if rows of pattern $P$ may start a walk and $0$ otherwise.
\end{lemma}

\begin{proof}
If $\pi(r)=\pi(r')$ then the map sending the value of $r$ at a column to the value of $r'$
at that column is a well-defined injection between the letters used, and extends to a
permutation $\sigma$ of the alphabet with $r'=\sigma(r)$. The edge condition is a statement
about equalities among entries, so $r\to s$ if and only if $\sigma(r)\to\sigma(s)$; and
$\pi(\sigma(s))=\pi(s)$. Hence $s\mapsto\sigma(s)$ is a bijection between the two successor
sets preserving patterns, which is the first claim. It says the chain is strongly lumpable
over $\pi$, so $\bigl(M_i^{m}\mathbf{1}\bigr)(r)$ depends on $r$ only through $\pi(r)$;
calling that value $f_m(\pi(r))$ gives the recursion, and summing over $r$ in each class
gives the displayed identity, the number of rows of pattern $P$ being the number of
injections of its $|P|$ classes into the alphabet.
\end{proof}

Assembling the $4$ lumped chains into one block-diagonal matrix $N$ of size $S=109$ and
putting the weights of Lemma~\ref{lem:inv} into the start vector,
$\mathbf{v}=\bigl(\binom{4}{i}D_{4-i}\varepsilon_P\bigr)_{i,P}$, gives
\begin{equation}\label{eq:walk}
4!\;a(n)\;=\;\mathbf{v}^{\!\top}N^{\,n}\mathbf{1} .
\end{equation}
The unlumped alphabet had $1300$ rows in all; the lumped chain has $S=109$ states.

\section{The criterion, and the computation}

Let $q(t)=t^{28}-\sum_i c_i\,t^{28-i}$ be the characteristic polynomial of
\eqref{eq:conj}.

\begin{lemma}\label{lem:crit}
For every $n$ with $m=n+1-1\ge28$,
\[
4!\Bigl(a(n)-\sum_i c_i\,a(n-i)\Bigr)\;=\;\mathbf{v}^{\!\top}N^{\,m-28}q(N)\mathbf{1} .
\]
Hence \eqref{eq:conj} holds from that point on if and only if
$\mathbf{v}^{\!\top}N^{j}q(N)\mathbf{1}=0$ for every $j\ge0$, and it is enough to check
$j=0,1,\dots,S-1$.
\end{lemma}

\begin{proof}
Substitute \eqref{eq:walk} into the left side and factor $N^{m-28}$ out of
$N^{m}-\sum_i c_iN^{m-i}$. The scalar $4!$ is nonzero, so it does not affect whether
the left side vanishes. For the last clause put $w=q(N)\mathbf{1}$: by Cayley--Hamilton
$N^{S}$ is an integer combination of $I,N,\dots,N^{S-1}$, so each
$\mathbf{v}^{\!\top}N^{j}w$ with $j\ge S$ is the corresponding combination of earlier
values.
\end{proof}

\begin{theorem}
The recurrence \eqref{eq:conj} holds for every $n>27$, which is the statement of
Section~1.
\end{theorem}

\begin{proof}
The lumped transition counts were obtained by enumerating, for one representative row of
each pattern, all admissible successors and sorting them by pattern --- the successors being
generated column by column, since the edge condition is a conjunction of one constraint per
column. The vector $w=q(N)\mathbf{1}$ was then formed by $28$ matrix--vector products
in exact integer arithmetic and $\mathbf{v}^{\!\top}N^{j}w$ evaluated for
$j=0,1,\dots$, again exactly. Every one of those integers vanishes from the index
corresponding to $n=27+1$ onwards, and the run continued until $w$ was the zero vector,
which settles all larger $j$ at once. Lemma~\ref{lem:crit} gives the theorem.
\end{proof}

\section{Verification}

Three checks, each independent of the algebra above.

First, the whole construction was compared against the entry's DATA: the right-hand side of
\eqref{eq:walk}, divided by $4!$, was evaluated at every one of the $12$
published indices and agrees exactly. This is a strong test of the modelling and not only of
the arithmetic --- the inversion of Lemma~\ref{lem:inv}, the reading of the local
condition, and the alignment of the entry's index with the number of rows would each
show up as a mismatch, and the quotient by $4!$ came out an integer at every index.

Second, the lumped chain was checked against the unlumped one: for the small shapes the walk
counts $\mathbf{s}_i^{\!\top}M_i^{m}\mathbf{1}$ were computed directly over the full
alphabet as well and agree with $\sum_P\varepsilon_Pf_m(P)$ term by term.

Third, the recurrence \eqref{eq:conj} was evaluated on the published terms in exact integer
arithmetic with no matrices involved, and the annihilation test of Lemma~\ref{lem:crit} was
run until the working vector was identically zero rather than to a sampled prefix.

\begin{thebibliography}{9}
\bibitem{oeis} The OEIS Foundation, \emph{The On-Line Encyclopedia of Integer
Sequences}, \url{https://oeis.org}, 2026. Sequence A210057.
\bibitem{stanley} R.~P.~Stanley, \emph{Enumerative Combinatorics}, Volume 1, 2nd ed.,
Cambridge University Press, 2011. (Transfer-matrix method, Section 4.7; Stirling and
falling-factorial inversion, Section 1.9.)
\bibitem{kemeny} J.~G.~Kemeny and J.~L.~Snell, \emph{Finite Markov Chains}, Springer,
1976. (Strong lumpability, Section 6.3.)
\bibitem{fs} P.~Flajolet and R.~Sedgewick, \emph{Analytic Combinatorics}, Cambridge
University Press, 2009.
\end{thebibliography}

\end{document}
